\documentclass[10pt,a4paper]{amsart}
\usepackage{amsmath,amssymb,amsthm,graphicx}
\usepackage{verbatim}
\newenvironment{reference}{\comment}{\endcomment}
%\newenvironment{reference}{}{}
\newenvironment{slogan}{\comment}{\endcomment}
\newenvironment{history}{\comment}{\endcomment}

% For commutative diagrams we use Xy-pic
\usepackage[all]{xy}

% We use 2cell for 2-commutative diagrams.
\xyoption{2cell}
% Theorem environments.
%
\theoremstyle{plain}
\newtheorem{theorem}[subsection]{Theorem}
\newtheorem{proposition}[subsection]{Proposition}
\newtheorem{lemma}[subsection]{Lemma}

\theoremstyle{definition}
\newtheorem{definition}[subsection]{Definition}
\newtheorem{example}[subsection]{Example}
\newtheorem{exercise}[subsection]{Exercise}
\newtheorem{situation}[subsection]{Situation}

\theoremstyle{remark}
\newtheorem{remark}[subsection]{Remark}
\newtheorem{remarks}[subsection]{Remarks}

\numberwithin{equation}{subsection}

% Macros
%
\def\lim{\mathop{\mathrm{lim}}\nolimits}
\def\colim{\mathop{\mathrm{colim}}\nolimits}
\def\Spec{\mathop{\mathrm{Spec}}}
\def\Hom{\mathop{\mathrm{Hom}}\nolimits}
\def\Ext{\mathop{\mathrm{Ext}}\nolimits}
\def\SheafHom{\mathop{\mathcal{H}\!\mathit{om}}\nolimits}
\def\SheafExt{\mathop{\mathcal{E}\!\mathit{xt}}\nolimits}
\def\Sch{\mathit{Sch}}
\def\Mor{\mathop{\mathrm{Mor}}\nolimits}
\def\Ob{\mathop{\mathrm{Ob}}\nolimits}
\def\Sh{\mathop{\mathit{Sh}}\nolimits}
\def\NL{\mathop{N\!L}\nolimits}
\def\CH{\mathop{\mathrm{CH}}\nolimits}
\def\proetale{{pro\text{-}\acute{e}tale}}
\def\etale{{\acute{e}tale}}
\def\QCoh{\mathit{QCoh}}
\def\Ker{\mathop{\mathrm{Ker}}}
\def\Im{\mathop{\mathrm{Im}}}
\def\Coker{\mathop{\mathrm{Coker}}}
\def\Coim{\mathop{\mathrm{Coim}}}

% Boxtimes
%
\DeclareMathSymbol{\boxtimes}{\mathbin}{AMSa}{"02}

%
% Macros for moduli stacks/spaces
%
\def\QCohstack{\mathcal{QC}\!\mathit{oh}}
\def\Cohstack{\mathcal{C}\!\mathit{oh}}
\def\Spacesstack{\mathcal{S}\!\mathit{paces}}
\def\Quotfunctor{\mathrm{Quot}}
\def\Hilbfunctor{\mathrm{Hilb}}
\def\Curvesstack{\mathcal{C}\!\mathit{urves}}
\def\Polarizedstack{\mathcal{P}\!\mathit{olarized}}
\def\Complexesstack{\mathcal{C}\!\mathit{omplexes}}
% \Pic is the operator that assigns to X its picard group, usage \Pic(X)
% \Picardstack_{X/B} denotes the Picard stack of X over B
% \Picardfunctor_{X/B} denotes the Picard functor of X over B
\def\Pic{\mathop{\mathrm{Pic}}\nolimits}
\def\Picardstack{\mathcal{P}\!\mathit{ic}}
\def\Picardfunctor{\mathrm{Pic}}
\def\Deformationcategory{\mathcal{D}\!\mathit{ef}}
\usepackage{cleveref}
\crefname{originallemma}{original lemma}{original lemmas}
\Crefname{originallemma}{Original lemma}{Original lemmas}
\makeatletter
\newlabel{lemma-associated-primes-localize}{{0310}{source}{}{sourcepointer-1}{}}
\newlabel{lemma-associated-primes-localize@cref}{{[originallemma][0][]0310}{[1][0][]source}{}{}{}}
\makeatother
\crefname{originallemma}{original lemma}{original lemmas}
\Crefname{originallemma}{Original lemma}{Original lemmas}
\makeatletter
\newlabel{lemma-ideal-nonzerodivisor}{{00LL}{source}{}{sourcepointer-2}{}}
\newlabel{lemma-ideal-nonzerodivisor@cref}{{[originallemma][0][]00LL}{[1][0][]source}{}{}{}}
\makeatother
\crefname{originallemma}{original lemma}{original lemmas}
\Crefname{originallemma}{Original lemma}{Original lemmas}
\makeatletter
\newlabel{proposition-minimal-primes-associated-primes}{{02CE}{source}{}{sourcepointer-3}{}}
\newlabel{proposition-minimal-primes-associated-primes@cref}{{[originallemma][0][]02CE}{[1][0][]source}{}{}{}}
\makeatother
\title{Having no embedded associated prime is equivalent to Serre's condition S one}
\author{The Stacks Project Authors}
\date{stacks-project:algebra.tex:local-snapshot-20261012; GNU FDL 1.2 or later}
\begin{document}
\maketitle
\noindent\textit{Editorial note.} All delivered mathematical text is the author's own, copied from the pinned Stacks Project chapter file algebra.tex; no statement or proof step was rewritten, reordered or removed.\par\medskip
\noindent\textit{Editorial note.} Preamble provenance, as required for this source: the frozen snapshot stages the chapter's own preamble file beside it (preamble.tex, 217 lines, 8116 bytes, the same shared Stacks-wide preamble the published chow and etale-cohomology cards use); the first line of the pinned chapter file is \textbackslash{}input\{preamble\}, so no preamble text lives inside the chapter itself. The card takes that file's lines 8-18 (the verbatim package and the three \textbackslash{}comment-based environments reference/slogan/history at lines 9-12, section 13.2, together with the Xy-pic setup) and its lines 154 to the runtime end of the file (the native theorem declarations and mathematical macros). Not one character was changed, no package was added, no macro was defined, and the excerpt serves only as preamble, never as a span coordinate.\par\medskip
\noindent\textit{Editorial note.} The same-chapter pointers are located at their \textbackslash{}label lines inside the pinned chapter file; the cross-chapter pointers are located in the project's own tag list staged in the snapshot (sources/stacks-chow/tags.txt), at the verbatim lines carrying their canonical tags.\par\medskip
\noindent\textit{Editorial note.} The counted proof refers to same-chapter propositions by label; they are recorded as navigation stubs pointed at their \textbackslash{}label lines, and the arguments consuming them stay inside the counted span.\par\medskip
\noindent\textit{Editorial note.} The counted spans contain no \textbackslash{}cite, so the card carries no bibliography.\par\medskip
\noindent\textit{Allowed standard prerequisites (editorial).}\begin{itemize}
\item Associated primes and their supports. Used throughout the counted proof (lines 44125-44145). Relates an embedded associated prime to a strictly smaller associated prime whose support contributes to the dimension.
\item Minimal primes of the support are associated. Invoked at the end of the counted proof (lines 44146-44148). Produces an associated prime inside a given one, which yields an embedded prime when the local depth condition fails.
\end{itemize}\medskip
\ifdefined\hypertarget\hypertarget{sourcepointer-1}{}\fi
\noindent\textit{Source reference.} Original lemma 0310: sources/stacks-algebra/pinned-original-source.tex, source lines 15611--15611. Navigation pointer for \textbackslash{}ref\{lemma-associated-primes-localize\} used inside the counted span (line 44131); the argument that consumes it is retained in the span.\par\medskip
\ifdefined\hypertarget\hypertarget{sourcepointer-2}{}\fi
\noindent\textit{Source reference.} Original lemma 00LL: sources/stacks-algebra/pinned-original-source.tex, source lines 15688--15688. Navigation pointer for \textbackslash{}ref\{lemma-ideal-nonzerodivisor\} used inside the counted span (line 44133); the argument that consumes it is retained in the span.\par\medskip
\ifdefined\hypertarget\hypertarget{sourcepointer-3}{}\fi
\noindent\textit{Source reference.} Original lemma 02CE: sources/stacks-algebra/pinned-original-source.tex, source lines 15405--15405. Navigation pointer for \textbackslash{}ref\{proposition-minimal-primes-associated-primes\} used inside the counted span (line 44146); the argument that consumes it is retained in the span.\par\medskip
\begin{lemma}
\label{lemma-criterion-no-embedded-primes}
Let $R$ be a Noetherian ring.
Let $M$ be a finite $R$-module.
The following are equivalent:
\begin{enumerate}
\item $M$ has no embedded associated prime, and
\item $M$ has property $(S_1)$.
\end{enumerate}
\end{lemma}

\begin{proof}
Let $\mathfrak p$ be an embedded associated prime of $M$.
Then there exists another associated prime $\mathfrak q$ of $M$
such that $\mathfrak p \supset \mathfrak q$. In particular this
implies that $\dim(\text{Supp}(M_{\mathfrak p})) \geq 1$ (since $\mathfrak q$
is in the support as well). On the other hand $\mathfrak pR_{\mathfrak p}$
is associated to $M_{\mathfrak p}$
(Lemma \ref{lemma-associated-primes-localize}) and hence
$\text{depth}(M_{\mathfrak p}) = 0$
(see Lemma \ref{lemma-ideal-nonzerodivisor}).
In other words $(S_1)$ does not hold.
Conversely, if $(S_1)$ does not hold then there exists a prime
$\mathfrak p$ such that $\dim(\text{Supp}(M_{\mathfrak p})) \geq 1$
and $\text{depth}(M_{\mathfrak p}) = 0$. Since
$\text{depth}(M_{\mathfrak p}) = 0$, we see that
$\mathfrak p \in \text{Ass}(M)$ by the two Lemmas
\ref{lemma-associated-primes-localize} and \ref{lemma-ideal-nonzerodivisor}.
Since $\dim(\text{Supp}(M_{\mathfrak p})) \geq 1$, there
is a prime $\mathfrak q \in \text{Supp}(M)$ with
$\mathfrak q \subset \mathfrak p$, $\mathfrak q \not = \mathfrak p$.
We can take such a $\mathfrak q$ that is minimal in
$\text{Supp}(M)$. Then by
Proposition \ref{proposition-minimal-primes-associated-primes}
we have $\mathfrak q \in \text{Ass}(M)$ and hence
$\mathfrak p$ is an embedded associated prime.
\end{proof}

\end{document}
